\section{符号说明}

本文使用的主要符号如表~\ref{tab:symbols} 所示。对于同时具有输入参数和派生量的符号，其具体属性在表中一并注明；功率变量均以区域、小时为索引，能量累计时采用时段长度 $\Delta t$ 进行换算。

{\scriptsize
\renewcommand{\arraystretch}{0.88}
\setlength{\tabcolsep}{3pt}
\begin{longtable}{lll}
\caption{主要符号及其含义}
\label{tab:symbols}\\
\toprule
符号 & 含义 & 类型或单位 \\
\midrule
\endfirsthead
\toprule
符号 & 含义 & 类型或单位 \\
\midrule
\endhead
\bottomrule
\endfoot
$i,r,t,s,k$ & 任务、区域、小时、开工时刻和任务类型索引 & 索引 \\
$\mathcal R$ & 六个计算区域构成的集合 & 集合 \\
$\mathcal K$ & 三类异构 AI 任务构成的集合 & 集合 \\
$\mathcal T_A$ & 任务到达主时域 & $\{0,\ldots,2399\}$ \\
$\mathcal T_C$ & 末端任务结清时域 & $\{2400,\ldots,2405\}$ \\
$\mathcal T_E$ & 电力与储能结算时域 & $\{0,\ldots,2406\}$ \\
$\mathcal S$ & 合法整数开工时刻集合 & $\{0,\ldots,2405\}$ \\
$\Delta t$ & 单个电力结算时段长度 & h \\
$a_i$ & 任务 $i$ 的到达小时 & h \\
$d_i,h_i$ & 任务持续分钟数及换算后的持续小时数 & min，h \\
$g_i$ & 任务 $i$ 持续占用的等效 GPU 数 & GPU \\
$k_i,o_i$ & 任务 $i$ 的类型和来源区域 & 类别，区域 \\
$F_i$ & 任务 $i$ 的有效完成截止时点 & h \\
$W_i^{\max}$ & 任务 $i$ 的最大允许等待时间 & h \\
$\ell_i^{\max}$ & 任务 $i$ 允许的最大网络时延 & ms \\
$\ell_{or}$ & 来源区域 $o$ 至执行区域 $r$ 的单向时延 & ms \\
$\omega_{it}(s)$ & 任务 $i$ 在开工时刻 $s$ 下与小时 $t$ 的重叠时长 & h \\
$\chi_{it}(s)$ & 任务 $i$ 在小时 $t$ 是否存在瞬时占用 & 二元参数 \\
$\Omega_i$ & 任务 $i$ 的合法区域--开工候选域 & 有限集合 \\
$x_{irs}$ & 任务 $i$ 是否在区域 $r$ 于时刻 $s$ 开工 & 二元变量 \\
$s_i,c_i,r_i$ & 任务实际开工时刻、完成时刻和执行区域 & h，h，区域 \\
$D_{rkt}$ & 区域 $r$ 类型 $k$ 在小时 $t$ 的实际到达 GPU 需求 & GPU \\
$D'_{rkt}$ & 用于滞后回归的历史需求输入 & GPU \\
$\widehat D_{rkt}$ & 区域--任务类型逐小时 GPU 需求预测值 & GPU \\
$u_{rkt}$ & Huber 回归截断前的线性预测量 & GPU \\
$z_t$ & 小时 $t$ 的日周期与周周期特征向量 & 无量纲 \\
$\beta,\gamma,\alpha_r,\delta_k$ & 回归系数、周期系数、区域效应和类型效应 & 参数 \\
$\lambda$ & Huber 回归的正则化系数 & 无量纲 \\
$\sigma_{rk}$ & 区域--任务类型序列的残差尺度 & GPU \\
$\rho_{\delta}(\cdot)$ & 阈值为 $\delta$ 的 Huber 损失函数 & 无量纲 \\
$G_r^{\max}$ & 区域 $r$ 的可调度 GPU 容量 & GPU \\
$P_r^{IT,\max}$ & 区域 $r$ 的最大 IT 功率 & MW \\
$P_r^{F,\max}$ & 区域 $r$ 的最大设施功率 & MW \\
$p_k^{GPU}$ & 类型 $k$ 的单位等效 GPU 平均 IT 功率 & MW/GPU \\
$\pi_r$ & 区域 $r$ 的能源使用效率 PUE & 无量纲 \\
$L_{rt}^{N}$ & 区域 $r$ 的不可迁移非 AI IT 负荷 & MW \\
$L_{rt}^{BAI}$ & 问题三固定的基准 AI IT 负荷 & MW \\
$L_{rt}^{AI,cap}$ & AI 任务形成的瞬时 IT 功率 & MW \\
$L_{rt}^{AI,avg}$ & AI 任务形成的小时平均 IT 负荷 & MW \\
$E_{rt}^{AI}$ & AI 任务在小时 $t$ 消耗的电量 & MWh \\
$L_{rt}^{F,avg}$ & 区域设施侧小时平均负荷 & MW \\
$U_{rt}^{peak}$ & 区域小时峰值 GPU 利用率 & 无量纲 \\
$U_{rt}^{avg}$ & 区域小时 GPU-hour 平均利用率 & 无量纲 \\
$R_{rt}^{av}$ & 区域可用新能源功率 & MW \\
$R_{rt}^{use}$ & 直接满足设施负荷的新能源功率 & MW \\
$R_{rt}^{ch}$ & 用于储能充电的新能源功率 & MW \\
$R_{rt}^{sell}$ & 新能源外送功率 & MW \\
$R_{rt}^{curt}$ & 弃置的新能源功率 & MW \\
$G_{rt}^{load}$ & 电网直接供给设施负荷的功率 & MW \\
$G_{rt}^{ch}$ & 电网供给储能充电的功率 & MW \\
$G_{rt}^{buy}$ & 区域从电网购入的总功率 & MW \\
$G_{rt}^{sell}$ & 区域向系统边界外送的功率 & MW \\
$I_r^{\max}$ & 区域最大购电功率 & MW \\
$X_r^{grid}$ & 区域电网侧最大外送功率 & MW \\
$X_r^{storage}$ & 储能信息限定的售电功率上限 & MW \\
$y_{rt}^{grid}$ & 区域购电与售电互斥状态 & 二元变量 \\
$C_{rt}$ & 储能总充电功率 & MW \\
$D_{rt}$ & 储能放电功率 & MW \\
$S_{rt}$ & 区域储能在小时 $t$ 运行后的 SOC & MWh \\
$B_r$ & 区域储能额定容量 & MWh \\
$S_r^0$ & 第 $0$ 小时运行前的初始 SOC & MWh \\
$S_r^{min},S_r^{max}$ & 区域储能 SOC 下限与上限 & MWh \\
$C_r^{\max},D_r^{\max}$ & 区域最大充电与放电功率 & MW \\
$\eta_r^c,\eta_r^d$ & 区域储能充电效率与放电效率 & 无量纲 \\
$N_{rt}$ & 区域净购电功率 $G_{rt}^{buy}-G_{rt}^{sell}$ & MW \\
$P_r^{peak}$ & 区域时域内非负峰值净购电功率 & MW \\
$V_r$ & 区域净购电绝对爬坡量 & MW \\
$C_{\mathrm{op}}$ & 购电支出扣除售电收入后的运行成本 & 元 \\
$E_{\mathrm{CO2}}$ & 电网购电产生的累计碳排放 & tCO$_2$ \\
$L_{\mathrm{net}}$ & 按 GPU 需求与执行时长加权的网络时延 & ms \\
$W_{\mathrm{net}}$ & 网络时延指标的有效权重和 & GPU$\cdot$h \\
$U_{\mathrm{RE}}$ & 新能源利用率 & 无量纲 \\
$W_{\mathrm{RE}}$ & 新能源利用率的可用新能源总量 & MWh \\
$Q_{\mathrm{RT}}$ & 实时任务即时开工满足率 & 无量纲 \\
$Q_{\mathrm{SLA}}$ & 网络时延 SLA 满足率 & 无量纲 \\
$Q_{\mathrm{deadline}}$ & 任务按期完成率 & 无量纲 \\
$Q_{\mathrm{wait}}$ & 弹性任务等待质量指标 & 无量纲 \\
$Q_{\mathrm{service}}$ & 综合服务质量指标 & 无量纲 \\
$\Psi_m$ & 实际有限候选搜索使用的代理综合评分 & 无量纲 \\
$w_j$ & 有效评价维度的实际等权系数 & 无量纲 \\
$d_j$ & 实际尺度 $\max\{|a_{j,m_0}|,1\}$ & 对应指标单位 \\
$a_{jm}$ & 候选方案 $m$ 在指标 $j$ 上的取值 & 对应指标单位 \\
$m_0$ & 标准化评分使用的基准方案索引 & 索引 \\
$N_{\mathrm{back}}$ & 冲突回溯节点预算 & 节点 \\
$N_{\mathrm{move}}^{q2}$ & 问题二单任务候选评价预算 & 次 \\
$N_{\mathrm{move}}^{q3}$ & 问题三储能动作候选评价预算 & 次 \\
$N_{\mathrm{move,total}}^{q4}$ & 问题四基准与全部情景共享的评价预算 & 次 \\
$D_{\mathrm{search}}$ & 启发式搜索的内部截止时刻 & 时间点 \\
$\rho_C$ & 情景中的碳约束水平参数 & 无量纲 \\
$\delta_P$ & 峰谷电价差调整系数 & 无量纲 \\
$\delta_S$ & 售电价格调整系数 & 无量纲 \\
$\delta_R$ & 新能源出力水平调整系数 & 无量纲 \\
$\xi$ & 新能源波动情景中的固定种子随机扰动 & 无量纲 \\
$\kappa_{rt}$ & 区域逐时电网购电碳强度 & tCO$_2$/MWh \\
$c_{rt}^{buy},c_{rt}^{sell}$ & 区域逐时购电价格和售电价格 & 元/MWh \\
$\mathrm{EFC}_r$ & 区域储能的等效完整循环次数 & 次 \\
\end{longtable}
}

表~\ref{tab:symbols} 同时区分了瞬时容量变量与小时能量变量，其中 $\chi_{it}(s)$ 被用于 GPU、IT 功率和设施功率的峰值审计，$\omega_{it}(s)$ 则被用于 GPU-hour 与 AI 电量计算。该区分能够保留非整小时任务的真实资源占用，避免小时平均功率掩盖瞬时超限。

能源变量按来源和去向被进一步拆分，使新能源守恒、电网购电成本和储能充电来源能够分别核算。由此，新能源利用率不会把弃电计入有效利用量，碳排放也不会遗漏电网为储能充电所提供的电量。

评分符号仅作用于分母有效且在候选之间存在真实差异的指标。若 $W_{\mathrm{net}}$、$W_{\mathrm{RE}}$ 或服务质量分量的相应分母不为正，则该指标被标记为不可用并退出综合评分，从而避免无依据的除零处理。
